\section{Introduction}

On April 2, 2026, \textit{Nature} published a feature article reporting that 2--6\% of papers accepted at major computer science conferences in 2025 contained ``hallucinated'' citations---references that look plausible but point to papers that do not exist \citep{naddaf2026}. The authors discuss an exclusive analysis, conducted in collaboration with the commercial firm Grounded AI, that estimated that more than 110{,}000 of the 7 million scholarly publications from 2025 contain invalid references generated by artificial intelligence.

The article was thorough, well-sourced, and alarming. It was also, in several respects, ironic.

\textit{Nature} partnered with \emph{Grounded AI} to analyze over 4{,}000 publications using a proprietary tool called \emph{Veracity}. Neither the tool, the dataset, the risk-scoring model, nor the 20{,}000 synthetic papers used to calibrate it were made publicly available. The methodology is described at a high level---enough to understand the approach, not enough to reproduce it. For an article about threats to scientific integrity, the absence of reproducible methods is striking.

The proposed solutions are similarly constrained. The article describes several detection tools: Grounded AI's \emph{Veracity} (commercial), GPTZero's screening tool (commercial partnership with ICLR), Frontiers' in-house AI (institutional), and CNRS's \emph{bibCheck} (available only to CNRS scientists). The only open-source option mentioned---\emph{CheckIfExist} \citep{abbonato2026}---receives a single sentence. The implicit message is that detection requires either money or institutional affiliation. Yet the researchers most vulnerable to hallucinated citations---editors of small journals, independent reviewers, scholars at under-resourced institutions---are precisely the ones who lack both.

This framing obscures a crucial fact: the infrastructure to solve this problem already exists, for free. CrossRef indexes over 150 million DOIs \citep{crossref2026}. OpenAlex covers 250 million scholarly works \citep{openalex2026}. Semantic Scholar, arXiv, and DataCite are all free and public \citep{semanticscholar2026}. The \textit{Nature} article itself relies on checking references against these same databases. The expensive part was never the data---it was the assumption that you need AI to interpret it.

This project started the same week that article was published. The first author read it as the editor of a small mathematics education journal---exactly the kind of journal with no budget for \emph{Grounded AI}, no institutional screening tools, and no defense against the problem \textit{Nature} described. The question was immediate and practical: \emph{what can I build, right now, with what's freely available?}

The answer, it turned out, was quite a lot. A deterministic validation pipeline using only free public APIs achieves a 0\% false-positive rate on 391 real citations and 100\% detection on fabricated citations that carry DOIs. No machine learning, no API keys, no subscription fees. The tool is limited---it is nearly blind to fabricated citations \emph{without} DOIs (4\% detection)---but its limitations are honest and precisely characterized.

Adding AI analysis (Gemini 2.5 Flash, free tier) raises detection to 94\% on fakes without DOIs. But it also introduces a 72\% false-positive rate on real arXiv preprints. The AI makes the same error we had already found and fixed in the deterministic logic: it treats ``I cannot verify this citation'' as ``this citation is probably fake.'' Grounded AI's own analysis for \textit{Nature} shows the same pattern---22 of their 100 most-flagged publications turned out to be genuine.

This recurring error is the central finding of this paper. It is not primarily a technical problem. It is an epistemological one: the conflation of \emph{unverifiable} with \emph{fabricated}, of absence of evidence with evidence of absence. We encountered it in our own code, watched an AI reproduce it independently, and found it reported at industrial scale in the very article that motivated this work. Any detection system---whether rule-based, AI-powered, or human---must grapple with this distinction. Most do not.

The remainder of this paper describes the tool and its evaluation (Methods), presents experimental results across deterministic and AI-assisted detection tiers (Results), develops the epistemological argument and its practical consequences (Discussion), and identifies the limitations that bound our claims (Limitations and Future Work). The complete codebase, all test data, and all experimental results are publicly available at \url{https://github.com/OhioMathTeacher/Ohio-Journal-of-School-Mathematics}.
